\documentclass[aspectratio=1610,12pt,notheorems]{beamer}

\usepackage[utf8x]{inputenc} \usepackage[russian]{babel}
\usepackage{amsmath,amssymb,amsthm,mathtools}
\usepackage{graphicx,caption,subcaption}
\usepackage{hyperref,natbib}
\usepackage{tikz,xcolor,makecell}
\usepackage{algorithm,algpseudocode}

\usetikzlibrary{arrows,backgrounds,patterns,%
	matrix,shapes,fit,calc,shadows,plotmarks,snakes}

\theoremstyle{plain}
\newtheorem{theorem}{Теорема}
\newtheorem{lemma}[theorem]{Лемма}

\theoremstyle{definition}
\newtheorem{definition}{Определение}
\newtheorem{problem}{Задача}

\usetheme[height=0.97cm]{Rochester}
\usecolortheme{dolphin}

\definecolor{hard}{RGB}{40,40,128} % {140,50,50}
\definecolor{mnsgold}{RGB}{255,255,130} % {240,230,200}

\setbeamercolor{headline}{bg=hard,fg=mnsgold}
\setbeamercolor*{frametitle}{parent=headline}

\setbeamercolor{structure}{fg=hard}
\setbeamercolor{subsection in head/foot}{bg=white,fg=hard}
\setbeamercolor{section in head/foot}{bg=hard,fg=mnsgold}
\setbeamercolor{block title}{bg=hard,fg=mnsgold}

\setbeamertemplate{navigation symbols}{}

\def\mitem{\medskip\item}
\def\ps{\\ [0.65cm]} \linespread{1.16} \parskip=2.6mm
\def\usl#1#2{\begin{block}{#1} #2 \end{block} \medskip\pause}

\def\ll{\left(} \def\rr{\right)}
\def\lag{\left\langle} \def\rag{\right\rangle}

\definecolor{failpos}{RGB}{230,30,20}
\definecolor{initpos}{RGB}{30,20,220}
\definecolor{turna}{RGB}{100,240,110}
\definecolor{turnb}{RGB}{250,140,110}


\title[Methodic keys to MNS]
    {\bfseries Организация олимпиады \\
	для своих учеников}

\author[\ ]
	{Б. А. Золотов,\quad «Математика НОН-СТОП»\\ \vspace{0.3cm}
		{\small Фонд «Время Науки»}}

\institute[\ ]{\ }

\date{27 ноября 2020}

%%%%%%%%%%%%%%
%%%%%%%%%%%%%%

\begin{document}

\frame{\titlepage}

\begin{frame} \frametitle{Олимпиада «Математика НОН-СТОП»}
\begin{itemize}
	\item Олимпиада по математике, \medskip
	\item Для 4–8 классов, \medskip
	\item Приспособлена для участников, не имеющих опыта кружков, \medskip
	\item Отсутствует отборочный тур, \medskip
	\item Предусмотрено разделение вариантов на базовые и профильные, \medskip
	\item Стимулирует продуманно подбирать стратегию решения задач. \medskip
\end{itemize}
\end{frame}

\begin{frame} \frametitle{Нас поддерживают}
\begin{itemize}
	\item Академия постдипломного педагогического образования: информационная поддержка, привлечение участников, \medskip
	\item Фонд Президентских грантов: финансирование олимпиады и других проектов в 2020–2021
\end{itemize}

	\vspace{0.45cm}
	\begin{center}
		\includegraphics[scale=0.5]{fpg/fpg}
	\end{center}
\end{frame}

\begin{frame}
	\ \\
	\noindent В 2019 году олимпиада «Математика НОН-СТОП» проводится
	с использованием гранта Президента Российской Федерации \\
	на развитие гражданского общества, предоставленного \\
	{\bfseries Фондом президентских грантов.}
	
	\vspace{0.45cm}
	\begin{center}
		\includegraphics[scale=0.6]{fpg}
	\end{center}
\end{frame}

\begin{frame}\frametitle{История олимпиады}
\begin{itemize}
	\item 2010 — первая олимпиада;
        %\mitem 2011 — 418 участников
%	\mitem 2012 — три площадки;
%	\mitem 2015 — первый профильный вариант;
	\item 2016 — 400 участников пишут базовый вариант, 92 --- профильный;\\
        $\phantom{2016}$ — поддержка Фонда «Время Науки»;
	%\mitem 2017 — задачи для 4 класса;\\
	%$\phantom{2017}$ — автоматическое распределение участников;
	\item 2018 — 847 участников пишут базовый вариант, 128 --- профильный;\\
        $\phantom{2018}$ — включение в Перечень региональных олимпиад и конкурсов\\
	$\phantom{2018 — }$\quad интеллектуальной направленности;\\
	$\phantom{2018}$ — поддержка Фонда Президентских грантов,\\
	$\phantom{2018 — }$\quad Комитета по образованию СПб;\\
	%$\phantom{2018}$ — этот семинар;\\
	\item 2019 — выход сборника задач;\\
        $\phantom{2019}$ — площадки в Бердске (Новосибирская обл.) и Гомеле (Беларусь);\\
        $\phantom{2019}$ — число участников приближается к 2000.
\end{itemize}\end{frame}

\begin{frame}\frametitle{Статистика олимпиады}
\begin{itemize}
	\item 12 площадок (на 2021 год — 20 соглашений);\\
	\item количество участников --- около 2000;\\
	\item пять городов: Санкт-Петербург, Бердск (Новосибирская обл.), \\
		Нов. Уренгой (ЯНАО), Гатчина (ЛО), \\
		Гомель (Беларусь);\\
	\item две страны;\\
	\item проблемы с часовыми поясами.
\end{itemize}\end{frame}

\begin{frame}\frametitle{Не верьте им!}
	Основная цель олимпиады — избегать появления «сообщества профессионалов», которые умеют её решать. \ps
	То, что вам сейчас расскажут, — не руководство по подготовке и не классификация задач будущей олимпиады; \ps
	А набор наблюдений, касающихся заданий прошлых лет.
\end{frame}


\begin{frame} \frametitle{Концепция заданий олимпиады}
\begin{itemize}
	\item Невозможно успеть решить все задачи, \medskip
	\item Есть задачи как для совсем неопытных участников, \\ так и для занимавшихся в кружках. \medskip
	\item Каждая задача разделена на пункты A, B и C, берётся \\ максимум баллов. \medskip
	\item Профильные варианты: новая обстановка для учеников матшкол.
\end{itemize}
\end{frame}

\begin{frame} \frametitle{Простые задания}
\usl{2020-4-4B}{
	У Вани есть доски для паркета размером $20 \times 10$ сантиметров, их можно распиливать пополам. Как Ване покрыть этими досками пол квадратной комнаты $1\text{ метр}\times 1\text{ метр}$ так, чтобы не было швов длиной более 30 сантиметров ни в одном из направлений?
}

\begin{center} \vspace{-0.4cm} \tikz[scale=0.4]
{
   \foreach \x / \y in {
	1/2, 1/6, 1/10, 2/1, 2/5, 2/9, 3/4, 3/8, 4/3, 4/7, 5/2, 5/6, 5/10, 6/1, 6/5, 6/9, 7/4, 7/8,
	8/3, 8/7, 9/1, 9/2, 9/6, 9/10
   } 
   \draw[fill opacity=0.6] (\x, \y) rectangle (\x cm + 2 cm, \y cm + 1 cm);
   \foreach \x / \y in {
	1/4, 1/8, 2/3, 2/7, 3/2, 3/6, 4/1, 4/5, 4/9, 5/4, 5/8, 6/3, 6/7, 7/2, 7/6, 8/1, 8/5, 8/9, 9/4, 9/8, 10/3, 10/7
   } 
   \draw[fill opacity=0.6] (\x, \y) rectangle (\x cm + 1 cm, \y cm + 2 cm);
   \foreach \x / \y in {
	1/1, 1/3, 1/7, 3/10, 5/1, 7/10, 10/5, 10/9
   } 
   \draw[fill opacity=0.6] (\x, \y) rectangle (\x cm + 1 cm, \y cm + 1 cm);
   \draw[thick,fill opacity=0.6] (1, 1) rectangle (11 cm, 11 cm);
}  \end{center}
\end{frame}

\begin{frame} \frametitle{Сложные задания}
Выпишем все числа от одного до десяти — но не в привычном порядке возрастания, а в алфавитном порядке: восемь, два, девять, десять, один, пять, семь, три, четыре, шесть. \medskip

\usl{2020-6-4B}{
Числа от 1 до 10'000'000'000 (десять миллиардов) выписали в алфавитном порядке. Перечислите первые десять из них.
}

(1) 18 (2) 18 миллионов (3) 18 миллионов 18 (4) 18 миллионов 18 тысяч \\
(5) 18 миллионов 18 тысяч 18 (6) \ldots восемь (7) \ldots восемьдесят \\
(8) \ldots 88 (9) \ldots 82 (10) \ldots 89.
\end{frame}

\begin{frame} \frametitle{Профильные задания}
{\it Система високосных лет} для числа $t$ — это последовательность натуральных чисел $(a_0, a_2, a_3, \ldots, a_n)$ такая, что $a_{i+1}$ делится на $a_i$, а также
	$$\frac{1}{a_0} - \frac{1}{a_1} + \frac{1}{a_2} - \frac{1}{a_3} + \ldots
	     + (-1)^{n} \cdot \frac{1}{a_n}\ =\ t.$$ \smallskip

Какой могла бы быть система високосных лет, если бы длина года составляла 365.21875, 365.17, 365.33 дней? Для любого ли рационального числа существует система високосных лет?
\end{frame}


\begin{frame} \frametitle{Площадки олимпиады}
\begin{itemize}
	\item Организация олимпиады в новых регионах, \medskip
	\item Увеличение числа участников, которое можно принять, \medskip
	\item Обеспечение удобства: написать олимпиаду ближе к дому, \medskip
	\item Выживание в ковид, \medskip
	\item Распространение опыта проведения олимпиад, \medskip
	\item Популяризация учреждения образования, \medskip
	\item Площадка, в т. ч., может быть организована «только для своих».
\end{itemize}
\end{frame}

\begin{frame} \frametitle{Информационные ресурсы олимпиады}
\begin{itemize}
	\item \url{mathnonstop.ru} — официальный сайт олимпиады со всеми объявлениями, условиями задач, итогами, \medskip
	\item \url{rs.mathnonstop.ru} — регистрационная система, работа со списком участников, загрузка результатов проверки, создание регистрационных листов, наклеек, дипломов, \medskip
	\item \url{vk.com/timeforscience} — ВКонтакте.\medskip
\end{itemize}
\end{frame}

\begin{frame} \frametitle{До олимпиады}
\begin{itemize}
	\item Начало регистрации участников — середина января–февраль, \medskip
	\item Нам нужно знать количество, размер, номера аудиторий, \medskip
	\item Нужен аккаунт организатора в \url{rs.mathnonstop.ru}, \medskip
	\item Волонтёры: $2n+5$, \medskip
	\item Можем разослать брошюры для занятий со школьниками. \medskip
\end{itemize}
\end{frame}

\begin{frame} \frametitle{Начало олимпиады}
\begin{itemize}
	\item Рассадка участников есть заранее, стол регистрации не нужен, \medskip
	\item Навигация $>$ регистрация, \medskip
	\item Посещаемость отмечается на регистрационных листах, \\ это самый главный документ, \medskip
	\item На тетради участников клеим наклейки, чтобы \\ не разбирать почерк школьников. \medskip
\end{itemize}
\end{frame}

\begin{frame} \frametitle{Продолжительность олимпиады}
\begin{itemize}
\item 4, 5 класс: $2.5$ часа; \medskip
\item 6 класс: 2 часа 45 минут; \medskip
\item 7, 8 класс, базовый вариант: 3 часа; \medskip
\item 7, 8 класс, профильный вариант: $3.5$ часа.
\end{itemize}
\end{frame}

\begin{frame} \frametitle{Инварианты по окончанию олимпиады}
\begin{itemize}
	\item Каждый пришедший должен сдать тетрадь, хотя бы пустую, \\
	но с наклейкой, \medskip
	\item Непришедшим участникам ставится $-1$, тогда система \\
	игнорирует их и не включает в таблицу результатов. \medskip
\end{itemize}
\end{frame}

\begin{frame} \frametitle{Проверка работ}
\begin{itemize}
	\item Желательно, организовывая олимпиаду, назначать \\ или находить проверяющих, \medskip
	\item Решения и критерии рассылаются проверяющим \\ почти сразу после олимпиады, \medskip
	\item Таблицы проверки экспортируются и импортируются в
	     \url{rs.mathnonstop.ru} \medskip
	\item Избранные работы~— случайные, спорные, по установленным \\ номерам~— загружаются туда же. \medskip
\end{itemize}
\end{frame}

\begin{frame} \frametitle{Распределение результатов}
\begin{itemize}
	\item Работы и результаты их проверки необходимо проконтролировать и отсмотреть. Председатель жюри перепроверит случайные работы из загруженных, то же нужно сделать и проверяющим: просмотреть несколько проверенных работ из другого региона. \bigskip
	\item Необходимые баллы для дипломов устанавливаются централизованно, а вот баллы на похвальные отзывы могут устанавливаться площадками.
\end{itemize}
\end{frame}

\def\fitem#1#2{\textcolor{hard}{\small (#1)}~~#2 \medskip \\}

\begin{frame} \frametitle{Заключение}

\textcolor{hard}{\bf Основные этапы проведения олимпиады}

\fitem{1}{Определение числа аудиторий и волонтёров}
\fitem{2}{Распределение участников по аудиториям в соответствие с рассадкой}
\fitem{3}{Сверка регистрационных листов}
\fitem{4}{Проверка работ}
\fitem{5}{Перепроверка работ, распределение наград \vspace{10mm}}

\hrule
\begin{center}
	{\LARGE Спасибо за внимание!} \smallskip \\
	{\footnotesize \textcolor{hard}{($*$)\quad}
		Заявки на разбор задач с \url{mathnonstop.ru}
		\phantom{($*$)\quad}}
\end{center} \vspace{2.4mm}
\end{frame}

\end{document}
